\unitchapter{2}{1}{Discrete Structures \& Optimization}{p2-01}

\unitmeta{Expected questions: 10--12 · Target: 9+ · Type: numerical + concept}

\begin{sylbox}

\begin{itemize}
\tightlist
\item[$\square$]
  Mathematical logic: propositional \& predicate logic, propositional equivalences, normal forms, predicates \& quantifiers, nested quantifiers, rules of inference
\item[$\square$]
  Sets \& relations: set operations, representation \& properties of relations, equivalence relations, partially ordering
\item[$\square$]
  Counting, mathematical induction \& discrete probability: basics of counting, pigeonhole principle, permutations \& combinations, inclusion-exclusion, mathematical induction, probability, Bayes\textquotesingle{} theorem
\item[$\square$]
  Group theory: groups, subgroups, semigroups, product \& quotients of algebraic structures, isomorphism, homomorphism, automorphism, rings, integral domains, fields, applications of group theory
\item[$\square$]
  Graph theory: simple graph, multigraph, weighted graph, paths \& circuits, shortest paths, Eulerian \& Hamiltonian, planar graph, graph colouring, bipartite graphs, trees \& rooted trees, prefix codes, tree traversals, spanning trees \& cut-sets
\item[$\square$]
  Boolean algebra: Boolean functions \& representation, simplification
\item[$\square$]
  Optimization: linear programming (graphical, simplex, duality), integer programming, transportation \& assignment, PERT-CPM
\end{itemize}

\end{sylbox}

\needspace{226pt}

\hypertarget{p2-01--1-mathematical-logic}{%
\section{Mathematical Logic}\label{p2-01--1-mathematical-logic}}

Logic is the grammar of mathematical reasoning and of computing: every condition in a program and every
gate in a circuit is a logical connective. In the examination, logic questions test three skills --- reading
a truth table, transforming formulas by known equivalences, and translating English into symbols.

\needspace{160pt}

\hypertarget{p2-01--11-connectives--truth-table}{%
\subsection{Connectives \& Truth Table}\label{p2-01--11-connectives--truth-table}}

The only row that makes \(p \rightarrow q\) false is the one where a true premise leads to a false
conclusion. A conditional with a false antecedent is "vacuously" true, which surprises students but follows
from treating \(p \rightarrow q\) as a promise that is broken only when \(p\) holds and \(q\) fails.

\begingroup\small

\par\medskip\noindent\hfil\begin{tabular}{@{}
  >{\raggedright\arraybackslash}p{(\columnwidth - 14\tabcolsep) * \real{0.0417}}
  >{\raggedright\arraybackslash}p{(\columnwidth - 14\tabcolsep) * \real{0.0417}}
  >{\raggedright\arraybackslash}p{(\columnwidth - 14\tabcolsep) * \real{0.0535}}
  >{\raggedright\arraybackslash}p{(\columnwidth - 14\tabcolsep) * \real{0.0659}}
  >{\raggedright\arraybackslash}p{(\columnwidth - 14\tabcolsep) * \real{0.0659}}
  >{\raggedright\arraybackslash}p{(\columnwidth - 14\tabcolsep) * \real{0.0659}}
  >{\raggedright\arraybackslash}p{(\columnwidth - 14\tabcolsep) * \real{0.0659}}
  >{\raggedright\arraybackslash}p{(\columnwidth - 14\tabcolsep) * \real{0.0659}}@{}}
\toprule\noalign{}
\begin{minipage}[b]{\linewidth}\raggedright
\textbf{p}
\end{minipage} & \begin{minipage}[b]{\linewidth}\raggedright
\textbf{q}
\end{minipage} & \begin{minipage}[b]{\linewidth}\raggedright
\textbf{¬p}
\end{minipage} & \begin{minipage}[b]{\linewidth}\raggedright
\textbf{p∧q}
\end{minipage} & \begin{minipage}[b]{\linewidth}\raggedright
\textbf{p∨q}
\end{minipage} & \begin{minipage}[b]{\linewidth}\raggedright
\textbf{p→q}
\end{minipage} & \begin{minipage}[b]{\linewidth}\raggedright
\textbf{p↔q}
\end{minipage} & \begin{minipage}[b]{\linewidth}\raggedright
\textbf{p⊕q}
\end{minipage} \\
\midrule\noalign{}
T & T & F & T & T & \textbf{T} & T & F \\
T & F & F & F & T & \textbf{F} & F & T \\
F & T & T & F & T & \textbf{T} & F & T \\
F & F & T & F & F & \textbf{T} & T & F \\
\bottomrule\noalign{}
\end{tabular}\hfil\null\par\medskip


\endgroup

\textbf{p → q is false ONLY when p is true and q is false.}

\begin{center}
\renewcommand{\arraystretch}{1.35}
\begin{tabular}{l c l}
\toprule
\textbf{Statement} & \textbf{Form} & \textbf{Equivalent to} \\
\midrule
conditional & $p \rightarrow q$ & its contrapositive \\
converse & $q \rightarrow p$ & the inverse \\
inverse & $\neg p \rightarrow \neg q$ & the converse \\
contrapositive & $\neg q \rightarrow \neg p$ & the conditional $p \rightarrow q$ \\
\bottomrule
\end{tabular}
\end{center}


\needspace{95pt}

\hypertarget{p2-01--12-important-equivalences}{%
\subsection{Important Equivalences}\label{p2-01--12-important-equivalences}}

\begin{formulatable}
Implication as disjunction & p \rightarrow q \equiv \neg p \lor q \\
Biconditional & p \leftrightarrow q \equiv (p\rightarrow q)\land(q\rightarrow p) \\
De Morgan's laws & \neg(p\land q)\equiv \neg p\lor\neg q,\quad \neg(p\lor q)\equiv\neg p\land\neg q \\
Absorption & p\lor(p\land q)\equiv p,\quad p\land(p\lor q)\equiv p \\
Exportation & (p\land q)\rightarrow r \equiv p\rightarrow(q\rightarrow r) \\
Common antecedent & (p\rightarrow q)\land(p\rightarrow r)\equiv p\rightarrow(q\land r) \\
Common consequent & (p\rightarrow r)\land(q\rightarrow r)\equiv (p\lor q)\rightarrow r \\
\end{formulatable}


\begin{itemize}
\tightlist
\item
  \textbf{Tautology}: always true (p ∨ ¬p). \textbf{Contradiction}: always false. \textbf{Contingency}: neither.
\item
  Functionally complete sets: \{∧, ¬\}, \{∨, ¬\}, \{NAND\}, \{NOR\}, \{→, ¬\}, \{→, F\}.
\end{itemize}

\needspace{95pt}

\hypertarget{p2-01--13-normal-forms}{%
\subsection{Normal Forms}\label{p2-01--13-normal-forms}}

\begin{itemize}
\tightlist
\item
  \textbf{CNF} = AND of ORs (product of sums); \textbf{DNF} = OR of ANDs (sum of products).
\item
  Number of Boolean functions on n variables = \textbf{2\^{}(2ⁿ)}.
\end{itemize}

\needspace{160pt}

\hypertarget{p2-01--14-rules-of-inference}{%
\subsection{Rules of Inference}\label{p2-01--14-rules-of-inference}}

Rules of inference are valid argument patterns: whenever the premises are true, the conclusion must be.
Proofs, resolution in artificial intelligence and many examination questions are chains of these few rules.

\begingroup\small

\par\medskip\noindent\hfil\begin{tabular}{@{}
  >{\raggedright\arraybackslash}p{(\columnwidth - 2\tabcolsep) * \real{0.2344}}
  >{\raggedright\arraybackslash}p{(\columnwidth - 2\tabcolsep) * \real{0.1661}}@{}}
\toprule\noalign{}
\begin{minipage}[b]{\linewidth}\raggedright
\textbf{Rule}
\end{minipage} & \begin{minipage}[b]{\linewidth}\raggedright
\textbf{Form}
\end{minipage} \\
\midrule\noalign{}
Modus Ponens & p, p→q ⊢ q \\
Modus Tollens & ¬q, p→q ⊢ ¬p \\
Hypothetical syllogism & p→q, q→r ⊢ p→r \\
Disjunctive syllogism & p∨q, ¬p ⊢ q \\
Addition & p ⊢ p∨q \\
Simplification & p∧q ⊢ p \\
Resolution & p∨q, ¬p∨r ⊢ q∨r \\
Universal instantiation & ∀x P(x) ⊢ P(c) \\
Existential generalisation & P(c) ⊢ ∃x P(x) \\
\bottomrule\noalign{}
\end{tabular}\hfil\null\par\medskip


\endgroup

Fallacies: affirming the consequent (q, p→q ⊢ p ✗), denying the antecedent (¬p, p→q ⊢ ¬q ✗).

\needspace{125pt}

\hypertarget{p2-01--15-predicates--quantifiers}{%
\subsection{Predicates \& Quantifiers}\label{p2-01--15-predicates--quantifiers}}

Predicates make statements about objects, and quantifiers say how many objects satisfy them. The most
common error is in translation: a universal statement uses implication, an existential one uses
conjunction.

\begin{formulatable}
Negating quantifiers & \neg\forall x\,P(x)\equiv\exists x\,\neg P(x),\quad \neg\exists x\,P(x)\equiv\forall x\,\neg P(x) \\
$\forall$ distributes over $\land$; $\exists$ over $\lor$ & \forall x(P\land Q)\equiv\forall xP\land\forall xQ,\quad \exists x(P\lor Q)\equiv\exists xP\lor\exists xQ \\
Order of mixed quantifiers & \exists x\,\forall y\,P(x,y)\ \Rightarrow\ \forall y\,\exists x\,P(x,y)\quad(\text{not conversely}) \\
``All students are smart'' — pair $\forall$ with $\rightarrow$ & \forall x\,\bigl(S(x)\rightarrow M(x)\bigr) \\
``Some students are smart'' — pair $\exists$ with $\land$ & \exists x\,\bigl(S(x)\land M(x)\bigr) \\
\end{formulatable}


\needspace{125pt}

\hypertarget{p2-01--2-sets--relations}{%
\section{Sets \& Relations}\label{p2-01--2-sets--relations}}

A relation on a set \(A\) is simply a subset of \(A \times A\), so with \(n\) elements there are \(2^{n^2}\) of
them. The counting formulas in the table follow by asking, for each ordered pair, how freely it may be
included.

\begin{itemize}
\tightlist
\item
  \textbar A ∪ B\textbar{} = \textbar A\textbar{} + \textbar B\textbar{} − \textbar A ∩ B\textbar; \textbar P(A)\textbar{} = 2ⁿ; \textbar A × B\textbar{} = mn.
\item
  Number of relations on A (\textbar A\textbar{} = n) = \textbf{2\^{}(n²)}.
\end{itemize}

\begingroup\small

\par\medskip\noindent\hfil\begin{tabular}{@{}
  >{\raggedright\arraybackslash}p{(\columnwidth - 2\tabcolsep) * \real{0.2072}}
  >{\raggedright\arraybackslash}p{(\columnwidth - 2\tabcolsep) * \real{0.4117}}@{}}
\toprule\noalign{}
\begin{minipage}[b]{\linewidth}\raggedright
\textbf{Relation type}
\end{minipage} & \begin{minipage}[b]{\linewidth}\raggedright
\textbf{Count on n elements}
\end{minipage} \\
\midrule\noalign{}
Reflexive & 2\^{}(n² − n) \\
Irreflexive & 2\^{}(n² − n) \\
Symmetric & 2\^{}(n(n+1)/\allowbreak{}2) \\
Antisymmetric & 2ⁿ · 3\^{}(n(n−1)/\allowbreak{}2) \\
Asymmetric & 3\^{}(n(n−1)/\allowbreak{}2) \\
Reflexive \& symmetric & 2\^{}(n(n−1)/\allowbreak{}2) \\
Equivalence relations & Bell number Bₙ (1, 1, 2, 5, 15, 52, 203 \ldots) \\
Functions A→B ( & A \\
One-one (m ≤ n) & n!/(n−m)! \\
Onto & Σ (−1)ᵏ C(n,k)(n−k)ᵐ \\
Bijections (m = n) & n! \\
\bottomrule\noalign{}
\end{tabular}\hfil\null\par\medskip


\endgroup

\needspace{161pt}

\hypertarget{p2-01--equivalence-relation--reflexive--symmetric--transitive--partitions-the-set-into-equivalence-classes}{%
\subsection{\texorpdfstring{Equivalence relation = Reflexive + Symmetric + Transitive → partitions the set into \textbf{equivalence classes}.}{Equivalence relation = Reflexive + Symmetric + Transitive → partitions the set into equivalence classes.}}\label{p2-01--equivalence-relation--reflexive--symmetric--transitive--partitions-the-set-into-equivalence-classes}}

\needspace{125pt}

\hypertarget{p2-01--partial-order-poset--reflexive--antisymmetric--transitive}{%
\subsection{Partial order (POSET) = Reflexive + Antisymmetric + Transitive.}\label{p2-01--partial-order-poset--reflexive--antisymmetric--transitive}}

\textbf{Hasse diagram} of divisibility on \{1, 2, 3, 6, 12\}:

\begin{figure}[!htbp]
\centering
\begin{tikzpicture}[v/.style={circ, minimum size=8mm}]
  \node[v] (1) at (0,0) {1}; \node[v] (2) at (-1.1,1.3) {2}; \node[v] (3) at (1.1,1.3) {3};
  \node[v] (6) at (0,2.6) {6}; \node[v] (12) at (0,3.9) {12};
  \draw[thinline] (1) -- (2) -- (6) -- (12) (1) -- (3) -- (6);
  \node[lbl, anchor=west, text width=60mm, align=left] at (2.4,2.0)
    {$a$ is joined below $b$ when $a$ divides $b$ and nothing lies between them.\\[4pt]
     Reflexive loops and implied (transitive) edges such as $1$–$6$ are omitted.};
\end{tikzpicture}
\caption{Hasse diagram of $\{1,2,3,6,12\}$ under divisibility.}
\end{figure}


\begin{itemize}
\tightlist
\item
  \textbf{Lattice}: every pair has LUB (join, ∨) and GLB (meet, ∧).
\item
  \textbf{Bounded}: has 0 (least) and 1 (greatest). \textbf{Complemented}: every element has a complement. \textbf{Distributive}: no sub-lattice isomorphic to M₃ (diamond) or N₅ (pentagon).
\item
  \textbf{Boolean algebra} = complemented + distributive lattice. Divisor lattice Dₙ is Boolean iff n is square-free.
\item
  Total order (chain): every pair comparable. Well-ordered: every nonempty subset has least element.
\end{itemize}

\begin{figure}[!htbp]
\centering
\begin{adjustbox}{max width=\linewidth}
\begin{tikzpicture}[v/.style={circ, minimum size=7mm, font=\small}]
  \begin{scope}
    \node[v] (0) at (0,0) {0}; \node[v] (a) at (-1.3,1.3) {$a$}; \node[v] (b) at (0,1.3) {$b$}; \node[v] (c) at (1.3,1.3) {$c$}; \node[v] (1) at (0,2.6) {1};
    \draw[thinline] (0) -- (a) -- (1) (0) -- (b) -- (1) (0) -- (c) -- (1);
    \node[capt] at (0,-0.8) {$M_3$ — the diamond};
  \end{scope}
  \begin{scope}[shift={(5,0)}]
    \node[v] (0) at (0,0) {0}; \node[v] (b) at (-1.1,0.95) {$b$}; \node[v] (a) at (-1.1,1.9) {$a$}; \node[v] (c) at (1.1,1.3) {$c$}; \node[v] (1) at (0,2.9) {1};
    \draw[thinline] (0) -- (b) -- (a) -- (1) (0) -- (c) -- (1);
    \node[capt] at (0,-0.8) {$N_5$ — the pentagon};
  \end{scope}
\end{tikzpicture}
\end{adjustbox}
\caption{The two smallest non-distributive lattices. A lattice is distributive exactly when it contains
neither as a sublattice.}
\end{figure}


\needspace{125pt}

\hypertarget{p2-01--3-counting--probability}{%
\section{Counting \& Probability}\label{p2-01--3-counting--probability}}

Counting questions reduce to one of a few models: ordered or unordered selection, with or without
repetition, in a line or a circle. Identify the model before reaching for a formula.

\begin{formulatable}
Permutations, combinations & {}^nP_r = \frac{n!}{(n-r)!},\qquad \binom{n}{r} = \frac{n!}{r!\,(n-r)!} \\
Circular arrangements; necklace or garland & (n-1)!\qquad\text{and}\qquad \frac{(n-1)!}{2} \\
Arrangements with repeated objects & \frac{n!}{p!\,q!\,r!} \\
Choosing $r$ from $n$ types with repetition (stars and bars) & \binom{n+r-1}{r} \\
Derangements & D_n = n!\sum_{k=0}^{n}\frac{(-1)^k}{k!};\quad D_1..D_5 = 0,1,2,9,44 \\
Catalan numbers (BSTs, bracketings) & C_n = \frac{1}{n+1}\binom{2n}{n} = 1,1,2,5,14,42,\dots \\
\end{formulatable}


\begin{itemize}
\tightlist
\item
  \textbf{Pigeonhole principle}: n items into m boxes ⇒ some box has at least ⌈n/m⌉ items.
  \emph{e.g. Min. people to ensure 2 share a birth month = 13.}
\item
  \textbf{Inclusion--exclusion}: \textbar A∪B∪C\textbar{} = ΣA − Σ(A∩B) + (A∩B∩C).
\item
  \textbf{Mathematical induction}: base case + (P(k) ⇒ P(k+1)). Strong induction assumes all P(1)\ldots P(k).
\end{itemize}

\needspace{95pt}

\hypertarget{p2-01--recurrence-relations}{%
\subsection{Recurrence relations}\label{p2-01--recurrence-relations}}

\begin{itemize}
\tightlist
\item
  Linear homogeneous: aₙ = c₁aₙ₋₁ + c₂aₙ₋₂ → characteristic equation r² − c₁r − c₂ = 0.
  Distinct roots r₁, r₂: aₙ = α r₁ⁿ + β r₂ⁿ; repeated root r: aₙ = (α + βn) rⁿ.
\item
  Fibonacci: Fₙ = Fₙ₋₁ + Fₙ₋₂. Tower of Hanoi: Tₙ = 2Tₙ₋₁ + 1 = 2ⁿ − 1.
\end{itemize}

\needspace{95pt}

\hypertarget{p2-01--probability}{%
\subsection{Probability}\label{p2-01--probability}}

\begin{formulatable}
Addition rule & P(A\cup B) = P(A)+P(B)-P(A\cap B) \\
Conditional probability & P(A\mid B) = \frac{P(A\cap B)}{P(B)} \\
Independence & P(A\cap B) = P(A)\,P(B) \\
Bayes' theorem & P(A\mid B) = \frac{P(B\mid A)\,P(A)}{P(B)} \\
Binomial distribution & P(X=k)=\binom{n}{k}p^k(1-p)^{n-k};\ \ \mu = np,\ \sigma^2 = npq \\
Poisson distribution & P(X=k) = \frac{e^{-\lambda}\lambda^k}{k!};\ \ \mu = \sigma^2 = \lambda \\
Expectation and variance & E[aX+b]=aE[X]+b,\quad \operatorname{Var}X = E[X^2]-(E[X])^2 \\
\end{formulatable}


\needspace{125pt}

\hypertarget{p2-01--4-algebraic-structures}{%
\section{Algebraic Structures}\label{p2-01--4-algebraic-structures}}

Abstract algebra studies sets with operations by listing the laws they obey. The same names recur in
computing: groups in cryptography, fields in error-correcting codes, Boolean algebras in circuit design.

\begin{figure}[!htbp]
\centering
\begin{tikzpicture}
  \foreach \w/\col/\name/\prop [count=\k] in {13.6/fgrey/Groupoid (magma)/closure, 11.4/fslate/Semigroup/{+ associativity}, 9.2/fsage/Monoid/{+ identity}, 7.0/fsand/Group/{+ inverses}, 4.8/frose/Abelian group/{+ commutativity}} {
    \pgfmathsetmacro{\yb}{0.28*(\k-1)}
    \pgfmathsetmacro{\yt}{5.0-0.62*(\k-1)}
    \filldraw[fill=\col, draw=fgink, line width=0.5pt, rounded corners=4pt] ({-\w/2},\yb) rectangle ({\w/2},\yt);
    \ifnum\k<5
      \node[font=\small, anchor=north] at (0,{\yt-0.08}) {\textbf{\name}\ \ \textit{\footnotesize\prop}};
    \else
      \node[font=\small] at (0,{(\yb+\yt)/2}) {\textbf{\name}\ \ \textit{\footnotesize\prop}};
    \fi
  }
\end{tikzpicture}
\caption{Algebraic structures with one operation, each adding one axiom to the one outside it.
Every abelian group is a group, every group a monoid, and so on.}
\end{figure}


\begin{figure}[!htbp]
\centering
\begin{adjustbox}{max width=\linewidth}
\begin{tikzpicture}[node distance=9mm]
  \node[slate, text width=28mm] (R) {\textbf{Ring}\\\scriptsize abelian group under $+$, semigroup under $\times$, distributive};
  \node[sage, text width=24mm, right=of R] (CR) {\textbf{Commutative ring}\\\scriptsize $ab = ba$};
  \node[sand, text width=27mm, right=of CR] (ID) {\textbf{Integral domain}\\\scriptsize unity, no zero divisors};
  \node[rose, text width=27mm, right=of ID] (F) {\textbf{Field}\\\scriptsize every $a \ne 0$ has an inverse};
  \draw[arr] (R) -- (CR); \draw[arr] (CR) -- (ID); \draw[arr] (ID) -- (F);
\end{tikzpicture}
\end{adjustbox}
\caption{Structures with two operations. $\mathbb{Z}$ is an integral domain but not a field; $\mathbb{Z}_n$ is a
field exactly when $n$ is prime; every finite integral domain is a field.}
\end{figure}


Key facts:

\begin{itemize}
\tightlist
\item
  \textbf{Lagrange\textquotesingle s theorem}: order of a subgroup divides the order of the group (finite).
\item
  Every group of \textbf{prime order} is cyclic and abelian. Every cyclic group is abelian.
\item
  Number of generators of cyclic group of order n = \textbf{φ(n)} (Euler totient).
\item
  Groups of order ≤ 5 are abelian. Smallest non-abelian group: S₃ (order 6).
\item
  (Zₙ, +ₙ) is a group; (Zₙ*, ×ₙ) group of units has φ(n) elements; Zₙ is a \textbf{field iff n is prime}.
\item
  Every finite integral domain is a field.
\item
  \textbf{Homomorphism}: f(a∗b) = f(a)∘f(b). \textbf{Isomorphism}: bijective homomorphism. \textbf{Automorphism}: isomorphism onto itself.
\item
  \textbf{Normal subgroup} H: gH = Hg ∀g; quotient group G/H has order \textbar G\textbar/\textbar H\textbar.
\item
  Order of an element a = smallest k with aᵏ = e.
\end{itemize}

\needspace{161pt}

\hypertarget{p2-01--5-graph-theory}{%
\section{Graph Theory}\label{p2-01--5-graph-theory}}

A graph is a set of vertices joined by edges --- an abstraction that models road networks, computer networks,
dependencies and scheduling conflicts. Most examination questions turn on a handful of counting results and
the characterisations of Euler graphs, planar graphs and bipartite graphs.

\needspace{95pt}

\hypertarget{p2-01--51-basic-results}{%
\subsection{Basic results}\label{p2-01--51-basic-results}}

\begin{formulatable}
Handshaking lemma & \sum_{v} \deg v = 2|E|\quad\Rightarrow\quad\text{odd-degree vertices are even in number} \\
Edges of $K_n$ (maximum for a simple graph) & \frac{n(n-1)}{2} \\
Edges of $K_{m,n}$ & mn \\
Tree on $n$ vertices; forest of $k$ trees & n-1;\qquad n-k \\
Labelled trees on $n$ vertices (Cayley) & n^{\,n-2}\ \ (= \text{spanning trees of } K_n) \\
Simple labelled graphs on $n$ vertices & 2^{\,n(n-1)/2} \\
\end{formulatable}


\needspace{95pt}

\hypertarget{p2-01--52-euler-vs-hamilton}{%
\subsection{Euler vs Hamilton}\label{p2-01--52-euler-vs-hamilton}}

\begin{figure}[!htbp]
\centering
\begin{adjustbox}{max width=\linewidth}
\begin{tikzpicture}[v/.style={circ, minimum size=6.5mm, font=\footnotesize}]
  \begin{scope}
    \node[v] (a) at (0,0) {a}; \node[v] (b) at (2,0) {b}; \node[v] (c) at (2,2) {c}; \node[v] (d) at (0,2) {d}; \node[v] (e) at (1,3.2) {e};
    \draw[fwine, line width=1.1pt] (a) -- (b) -- (c) -- (d) -- (a) (d) -- (e) -- (c);
    \node[lbl, align=center] at (1,-0.9) {all degrees even\\$\Rightarrow$ an Euler circuit exists};
    \node[capt] at (1,4.1) {Euler: every \emph{edge} once};
  \end{scope}
  \begin{scope}[shift={(6.2,0)}]
    \node[v] (p) at (0,0) {p}; \node[v] (q) at (2,0) {q}; \node[v] (r) at (2.6,1.6) {r}; \node[v] (s) at (1,2.7) {s}; \node[v] (t) at (-0.6,1.6) {t};
    \draw[thinline] (p) -- (r) (q) -- (s) (s) -- (p);
    \draw[fwine, line width=1.1pt] (p) -- (q) -- (r) -- (s) -- (t) -- (p);
    \node[lbl, align=center] at (1,-0.9) {highlighted cycle visits\\every vertex exactly once};
    \node[capt] at (1,4.1) {Hamilton: every \emph{vertex} once};
  \end{scope}
\end{tikzpicture}
\end{adjustbox}
\caption{Euler and Hamilton problems. Euler's condition is a simple degree test; deciding whether a
Hamiltonian cycle exists is NP-complete, though Dirac's and Ore's conditions guarantee one.}
\end{figure}


\begin{itemize}
\tightlist
\item
  Kₙ is Eulerian iff n is odd. Kₙ has (n−1)!/2 Hamiltonian cycles.
\end{itemize}

\needspace{125pt}

\hypertarget{p2-01--53-planar-graphs}{%
\subsection{Planar Graphs}\label{p2-01--53-planar-graphs}}

A graph is planar if it can be drawn in the plane without edges crossing. Euler\textquotesingle s formula links the
numbers of vertices, edges and faces of any such drawing, and the edge bounds derived from it give a quick
test for non-planarity.

\begin{formulatable}
Euler's formula (connected planar graph) & V - E + F = 2\qquad(k\text{ components: } V-E+F = k+1) \\
Simple planar graph, $V \ge 3$ & E \le 3V - 6 \\
Planar graph with no triangles & E \le 2V - 4 \\
Kuratowski's theorem & G\text{ planar} \iff \text{no subdivision of } K_5 \text{ or } K_{3,3} \\
\end{formulatable}


\begin{figure}[!htbp]
\centering
\begin{adjustbox}{max width=\linewidth}
\begin{tikzpicture}[v/.style={circ, minimum size=6mm, font=\footnotesize}]
  \begin{scope}
    \foreach \x/\n in {0/a,1.3/b,2.6/c} \node[v, fill=fslate] (\n) at (\x,1.8) {\n};
    \foreach \x/\n in {0/x,1.3/y,2.6/z} \node[v, fill=fsand] (\n) at (\x,0) {\n};
    \foreach \u in {a,b,c} \foreach \w in {x,y,z} \draw[thinline] (\u) -- (\w);
    \node[capt] at (1.3,-0.8) {$K_{3,3}$ — non-planar};
  \end{scope}
  \begin{scope}[shift={(5.2,0.9)}]
    \foreach \i in {0,...,4} \node[v, fill=frose] (k\i) at ({90+72*\i}:1.15) {};
    \foreach \i in {0,...,4} \foreach \j in {0,...,4} {\ifnum\i<\j \draw[thinline] (k\i) -- (k\j);\fi}
    \node[capt] at (0,-1.7) {$K_5$ — non-planar};
  \end{scope}
  \begin{scope}[shift={(9.6,0)}]
    \node[v, fill=fsage] (1) at (1.1,2.0) {1}; \node[v, fill=fsage] (2) at (0,0) {2}; \node[v, fill=fsage] (3) at (2.2,0) {3}; \node[v, fill=fsage] (4) at (1.1,0.75) {4};
    \draw[thinline] (1) -- (2) -- (3) -- (1) (4) -- (1) (4) -- (2) (4) -- (3);
    \node[capt] at (1.1,-0.8) {$K_4$ — planar};
  \end{scope}
\end{tikzpicture}
\end{adjustbox}
\caption{$K_{3,3}$ and $K_5$ cannot be drawn without crossings; $K_4$ can, with one vertex placed inside the triangle.}
\end{figure}


\needspace{125pt}

\hypertarget{p2-01--54-colouring}{%
\subsection{Colouring}\label{p2-01--54-colouring}}

Colouring assigns colours to vertices so that adjacent vertices differ --- the model behind timetabling and
register allocation. The chromatic number is easy to bound but hard to compute in general.

\begin{itemize}
\tightlist
\item
  \textbf{Chromatic number χ(G)}: min colours so adjacent vertices differ.
\item
  χ(Kₙ) = n; χ(bipartite with ≥1 edge) = 2; χ(tree, n ≥ 2) = 2; χ(cycle Cₙ) = 2 if n even, 3 if odd; χ(Wheel Wₙ) = 3 or 4.
\item
  \textbf{Four colour theorem}: planar ⇒ χ ≤ 4.
\item
  G is \textbf{bipartite ⇔ no odd cycle ⇔ 2-colourable}.
\item
  Chromatic polynomial of tree with n vertices: k(k−1)ⁿ⁻¹; of Kₙ: k(k−1)\ldots(k−n+1).
\end{itemize}

\needspace{95pt}

\hypertarget{p2-01--55-other-definitions}{%
\subsection{Other definitions}\label{p2-01--55-other-definitions}}

\begin{itemize}
\tightlist
\item
  \textbf{Independent set}, \textbf{clique}, \textbf{vertex cover}: α(G) + β(G) = n (independence number + vertex cover number).
\item
  \textbf{Matching}; perfect matching. Edge chromatic number (Vizing): Δ ≤ χ′ ≤ Δ+1.
\item
  \textbf{Cut vertex / articulation point}, \textbf{bridge}, \textbf{cut-set}. Vertex connectivity κ ≤ edge connectivity λ ≤ δ (min degree) --- Whitney.
\item
  \textbf{Isomorphism} invariants: same \#vertices, \#edges, degree sequence, cycles.
\item
  \textbf{Prefix codes}: Huffman coding gives optimal prefix-free code using a binary tree.
\end{itemize}

\needspace{95pt}

\hypertarget{p2-01--6-boolean-algebra}{%
\section{Boolean Algebra}\label{p2-01--6-boolean-algebra}}

\begin{formulatable}
Idempotent laws & x + x = x,\qquad x\cdot x = x \\
Absorption & x + xy = x,\qquad x(x+y) = x \\
Redundant literal & x + \bar{x}y = x + y \\
Distribution of $+$ over $\cdot$ & (x+y)(x+z) = x + yz \\
Consensus theorem & xy + \bar{x}z + yz = xy + \bar{x}z \\
Duality & \text{interchange } + \leftrightarrow \cdot \text{ and } 0 \leftrightarrow 1 \\
\end{formulatable}


K-maps and minimisation are covered in \protect\hyperlink{p2-02}{Unit 2}.

\needspace{161pt}

\hypertarget{p2-01--7-optimization}{%
\section{Optimization}\label{p2-01--7-optimization}}

Optimization chooses the best value of an objective subject to constraints. In linear programming both
objective and constraints are linear, so the feasible region is a convex polygon (or polyhedron) and the
optimum lies at a corner --- the insight on which the simplex method is built.

\needspace{95pt}

\hypertarget{p2-01--71-linear-programming-lpp}{%
\subsection{Linear Programming (LPP)}\label{p2-01--71-linear-programming-lpp}}

\begin{figure}[!htbp]
\centering
\begin{tikzpicture}[node distance=10mm]
  \node[slate, text width=34mm] (A) {\textbf{Formulate}\\\footnotesize variables, objective, constraints};
  \node[diamondbox, fill=fsand, right=of A] (B) {two\\variables?};
  \node[sage, text width=30mm, right=14mm of B, yshift=8mm] (C) {\textbf{Graphical method}\\\footnotesize test corner points};
  \node[sage, text width=30mm, right=14mm of B, yshift=-8mm] (D) {\textbf{Simplex method}};
  \node[rose, text width=26mm, right=of D] (E) {duality,\\dual simplex};
  \draw[arr] (A) -- (B); \draw[arr] (B) |- node[lbl, above, pos=0.75] {yes} (C); \draw[arr] (B) |- node[lbl, below, pos=0.75] {no} (D); \draw[arr] (D) -- (E);
\end{tikzpicture}
\caption{Choosing a method for a linear programme.}
\end{figure}


\textbf{Graphical example}: Max Z = 3x + 5y, s.t. x + 2y ≤ 8, 3x + 2y ≤ 12, x, y ≥ 0.

\begin{figure}[!htbp]
\centering
\begin{adjustbox}{max width=\linewidth}
\begin{tikzpicture}
\begin{axis}[width=10cm, height=8cm, axis lines=middle, xmin=0, xmax=8.6, ymin=0, ymax=6.8, xlabel={$x$}, ylabel={$y$},
    xtick={0,2,4,6,8}, ytick={0,2,4,6}, tick label style={font=\footnotesize}, clip=false, axis line style={-{Stealth[length=2mm]}}]
  \addplot[fill=fsage!80, draw=none] coordinates {(0,0) (4,0) (2,3) (0,4)} -- cycle;
  \addplot[fgink, line width=0.7pt, domain=0:8] {(8-x)/2} node[pos=0.88, above, font=\footnotesize, sloped] {$x+2y=8$};
  \addplot[fgink, line width=0.7pt, domain=0:4] {(12-3*x)/2} node[pos=0.12, right, font=\footnotesize] {$3x+2y=12$};
  \addplot[fwine, dashed, line width=0.7pt, domain=0:7] {(21-3*x)/5} node[pos=0.97, below, font=\footnotesize, text=fwine] {$Z=21$};
  \fill[fgink] (axis cs:0,0) circle (2pt) (axis cs:4,0) circle (2pt) (axis cs:2,3) circle (2pt) (axis cs:0,4) circle (2pt);
  \node[font=\footnotesize, anchor=south west] at (axis cs:4.05,0.08) {(4,0): 12};
  \node[font=\footnotesize, anchor=south west] at (axis cs:2.05,3.05) {(2,3): \textbf{21} — maximum};
  \node[font=\footnotesize, anchor=south west] at (axis cs:0.08,4.08) {(0,4): 20};
  \node[font=\small] at (axis cs:1.2,1.3) {feasible region};
\end{axis}
\end{tikzpicture}
\end{adjustbox}
\caption{Graphical solution of: maximise $Z = 3x+5y$ subject to $x+2y\le 8$, $3x+2y\le 12$, $x,y\ge 0$.
The objective line, moved outward, last touches the region at the corner $(2,3)$.}
\end{figure}


\begin{itemize}
\tightlist
\item
  Optimal solution occurs at a \textbf{corner point} (extreme point) of the convex feasible region.
\item
  Special cases: \textbf{unbounded}, \textbf{infeasible}, \textbf{alternate optima} (objective parallel to a binding constraint), \textbf{degeneracy} (basic variable = 0).
\item
  Slack (for ≤), surplus (for ≥), artificial variables (Big-M / Two-phase).
\item
  Simplex: entering variable = most negative Cⱼ − Zⱼ (max) ; leaving variable = \textbf{minimum ratio test}.
\item
  \textbf{Duality}: primal max ⇔ dual min; \#constraints ↔ \#variables; dual of dual = primal. Strong duality: optimal values equal. Weak duality: any feasible dual value bounds the primal.
\item
  \textbf{Dual simplex}: starts optimal but infeasible, keeps optimality till feasibility.
\end{itemize}

\needspace{125pt}

\hypertarget{p2-01--72-transportation-problem}{%
\subsection{Transportation Problem}\label{p2-01--72-transportation-problem}}

The transportation problem ships goods from sources to destinations at minimum cost. It is a linear
programme with special structure, solved in two stages: find a starting allocation, then improve it until
no cheaper route exists.

\begin{itemize}
\tightlist
\item
  m sources, n destinations; balanced if supply = demand (else add dummy).
\item
  A basic feasible solution has \textbf{m + n − 1} allocations (fewer → degenerate).
\item
  Initial solution: \textbf{North-West Corner}, \textbf{Least Cost}, \textbf{Vogel\textquotesingle s Approximation (VAM --- best)}.
\item
  Optimality: \textbf{MODI (u-v) method} or Stepping stone.
\end{itemize}

\needspace{125pt}

\hypertarget{p2-01--73-assignment-problem}{%
\subsection{Assignment Problem}\label{p2-01--73-assignment-problem}}

Assigning \(n\) jobs to \(n\) people one-to-one is a transportation problem with every supply and demand equal
to one; the Hungarian method exploits this to solve it by simple row and column reductions.

\begin{itemize}
\tightlist
\item
  n jobs to n persons, one-to-one; solved by \textbf{Hungarian method} (row reduction, column reduction, cover zeros with min lines; if lines = n ⇒ optimal).
\item
  Special case of transportation (all supplies/demands = 1); highly degenerate.
\end{itemize}

\needspace{95pt}

\hypertarget{p2-01--74-integer-programming}{%
\subsection{Integer Programming}\label{p2-01--74-integer-programming}}

\begin{itemize}
\tightlist
\item
  Variables restricted to integers: \textbf{Branch and Bound}, \textbf{Gomory\textquotesingle s cutting plane}.
\end{itemize}

\needspace{125pt}

\hypertarget{p2-01--75-pert-cpm}{%
\subsection{PERT-CPM}\label{p2-01--75-pert-cpm}}

Large projects are planned as networks of activities. The critical path is the longest path through the
network: any delay on it delays the whole project, whereas activities off the path have slack (float).

\begin{figure}[!htbp]
\centering
\begin{tikzpicture}[ev/.style={ink, circle split, fill=fpaper, minimum size=12mm, inner sep=1pt, font=\scriptsize,
    blur shadow={shadow blur steps=6, shadow xshift=0.7pt, shadow yshift=-0.9pt, shadow opacity=22}}]
  \node[ev] (1) at (0,0) {\textbf{1} \nodepart{lower} 0 | 0};
  \node[ev] (2) at (3,1.6) {\textbf{2} \nodepart{lower} 3 | 3};
  \node[ev] (3) at (3,-1.6) {\textbf{3} \nodepart{lower} 4 | 6};
  \node[ev] (4) at (6.4,0) {\textbf{4} \nodepart{lower} 8 | 8};
  \node[ev] (5) at (9.4,0) {\textbf{5} \nodepart{lower} 11 | 11};
  \draw[arr, fwine, line width=1.2pt] (1) -- node[edgelbl, above left] {A = 3} (2);
  \draw[arr] (1) -- node[edgelbl, below left] {B = 4} (3);
  \draw[arr, fwine, line width=1.2pt] (2) -- node[edgelbl, above right] {C = 5} (4);
  \draw[arr] (3) -- node[edgelbl, below right] {D = 2} (4);
  \draw[arr, fwine, line width=1.2pt] (4) -- node[edgelbl, above] {E = 3} (5);
  \node[lbl, anchor=west] at (10.3,0.4) {earliest | latest};
  \node[lbl, anchor=west, text=fwine] at (10.3,-0.2) {critical path};
\end{tikzpicture}
\caption{An activity-on-arrow network. Each event shows its earliest and latest times; where they are
equal the event is critical. The critical path A–C–E (11 days) fixes the project duration, and activities
B and D share 2 days of float.}
\end{figure}


Paths: A-C-E = 3+5+3 = \textbf{11} (critical), B-D-E = 4+2+3 = 9. Project duration = 11.

\begingroup\small

\par\medskip\noindent\hfil\begin{tabular}{@{}
  >{\raggedright\arraybackslash}p{(\columnwidth - 2\tabcolsep) * \real{0.2689}}
  >{\raggedright\arraybackslash}p{(\columnwidth - 2\tabcolsep) * \real{0.2806}}@{}}
\toprule\noalign{}
\begin{minipage}[b]{\linewidth}\raggedright
\textbf{Concept}
\end{minipage} & \begin{minipage}[b]{\linewidth}\raggedright
\textbf{Formula}
\end{minipage} \\
\midrule\noalign{}
Earliest start (forward pass) & ES = max(EF of predecessors) \\
Latest finish (backward pass) & LF = min(LS of successors) \\
Total float & LS − ES = LF − EF \\
Critical activity & Total float = 0 \\
PERT expected time & tₑ = (a + 4m + b)/6 \\
PERT variance & σ² = ((b − a)/6)² \\
Probability & Z = (T − Tₑ)/σ\_path \\
\bottomrule\noalign{}
\end{tabular}\hfil\null\par\medskip


\endgroup

\begin{itemize}
\tightlist
\item
  \textbf{CPM}: deterministic time, cost-oriented (crashing). \textbf{PERT}: probabilistic (β-distribution), event-oriented.
\item
  \textbf{Crashing}: reducing duration by adding resources; crash cost slope = (crash cost − normal cost)/(normal time − crash time).
\item
  \textbf{Resource levelling}: smooth resource usage without changing project duration (use floats).
\end{itemize}

\needspace{191pt}

\hypertarget{p2-01--8-deeper-dive--worked-simplex-transportation-assignment-groups--recurrences}{%
\section{Deeper Dive --- Worked Simplex, Transportation, Assignment, Groups \& Recurrences}\label{p2-01--8-deeper-dive--worked-simplex-transportation-assignment-groups--recurrences}}

\needspace{155pt}

\hypertarget{p2-01--81-simplex-method--full-worked-example}{%
\subsection{Simplex Method --- Full Worked Example}\label{p2-01--81-simplex-method--full-worked-example}}

The simplex method moves from corner to corner of the feasible region, always improving the objective, and
stops when no adjacent corner is better. Each tableau below is one such corner.

Max Z = 3x₁ + 5x₂ subject to x₁ ≤ 4, 2x₂ ≤ 12, 3x₁ + 2x₂ ≤ 18, x₁, x₂ ≥ 0.
Add slacks s₁, s₂, s₃.

\textbf{Tableau 0} (Z row written as Z − 3x₁ − 5x₂ = 0)

\begingroup\small

\par\medskip\noindent\hfil\begin{tabular}{@{}
  >{\raggedright\arraybackslash}p{(\columnwidth - 14\tabcolsep) * \real{0.0893}}
  >{\raggedright\arraybackslash}p{(\columnwidth - 14\tabcolsep) * \real{0.0535}}
  >{\raggedright\arraybackslash}p{(\columnwidth - 14\tabcolsep) * \real{0.0785}}
  >{\raggedright\arraybackslash}p{(\columnwidth - 14\tabcolsep) * \real{0.0535}}
  >{\raggedright\arraybackslash}p{(\columnwidth - 14\tabcolsep) * \real{0.0535}}
  >{\raggedright\arraybackslash}p{(\columnwidth - 14\tabcolsep) * \real{0.0535}}
  >{\raggedright\arraybackslash}p{(\columnwidth - 14\tabcolsep) * \real{0.0838}}
  >{\raggedright\arraybackslash}p{(\columnwidth - 14\tabcolsep) * \real{0.1396}}@{}}
\toprule\noalign{}
\begin{minipage}[b]{\linewidth}\raggedright
\textbf{Basic}
\end{minipage} & \begin{minipage}[b]{\linewidth}\raggedright
\textbf{x₁}
\end{minipage} & \begin{minipage}[b]{\linewidth}\raggedright
\textbf{x₂}
\end{minipage} & \begin{minipage}[b]{\linewidth}\raggedright
\textbf{s₁}
\end{minipage} & \begin{minipage}[b]{\linewidth}\raggedright
\textbf{s₂}
\end{minipage} & \begin{minipage}[b]{\linewidth}\raggedright
\textbf{s₃}
\end{minipage} & \begin{minipage}[b]{\linewidth}\raggedright
\textbf{RHS}
\end{minipage} & \begin{minipage}[b]{\linewidth}\raggedright
\textbf{Ratio}
\end{minipage} \\
\midrule\noalign{}
s₁ & 1 & 0 & 1 & 0 & 0 & 4 & --- \\
s₂ & 0 & \textbf{2} & 0 & 1 & 0 & 12 & 12/2 = \textbf{6} ← \\
s₃ & 3 & 2 & 0 & 0 & 1 & 18 & 18/2 = 9 \\
Z & −3 & \textbf{−5} ↑ & 0 & 0 & 0 & 0 & \\
\bottomrule\noalign{}
\end{tabular}\hfil\null\par\medskip


\endgroup

x₂ enters (most negative), s₂ leaves (minimum ratio). Pivot = 2.

\textbf{Tableau 1}

\begingroup\small

\par\medskip\noindent\hfil\begin{tabular}{@{}
  >{\raggedright\arraybackslash}p{(\columnwidth - 14\tabcolsep) * \real{0.0893}}
  >{\raggedright\arraybackslash}p{(\columnwidth - 14\tabcolsep) * \real{0.0785}}
  >{\raggedright\arraybackslash}p{(\columnwidth - 14\tabcolsep) * \real{0.0535}}
  >{\raggedright\arraybackslash}p{(\columnwidth - 14\tabcolsep) * \real{0.0535}}
  >{\raggedright\arraybackslash}p{(\columnwidth - 14\tabcolsep) * \real{0.0597}}
  >{\raggedright\arraybackslash}p{(\columnwidth - 14\tabcolsep) * \real{0.0535}}
  >{\raggedright\arraybackslash}p{(\columnwidth - 14\tabcolsep) * \real{0.0838}}
  >{\raggedright\arraybackslash}p{(\columnwidth - 14\tabcolsep) * \real{0.1257}}@{}}
\toprule\noalign{}
\begin{minipage}[b]{\linewidth}\raggedright
\textbf{Basic}
\end{minipage} & \begin{minipage}[b]{\linewidth}\raggedright
\textbf{x₁}
\end{minipage} & \begin{minipage}[b]{\linewidth}\raggedright
\textbf{x₂}
\end{minipage} & \begin{minipage}[b]{\linewidth}\raggedright
\textbf{s₁}
\end{minipage} & \begin{minipage}[b]{\linewidth}\raggedright
\textbf{s₂}
\end{minipage} & \begin{minipage}[b]{\linewidth}\raggedright
\textbf{s₃}
\end{minipage} & \begin{minipage}[b]{\linewidth}\raggedright
\textbf{RHS}
\end{minipage} & \begin{minipage}[b]{\linewidth}\raggedright
\textbf{Ratio}
\end{minipage} \\
\midrule\noalign{}
s₁ & 1 & 0 & 1 & 0 & 0 & 4 & 4 \\
x₂ & 0 & 1 & 0 & 1/2 & 0 & 6 & --- \\
s₃ & \textbf{3} & 0 & 0 & −1 & 1 & 6 & 6/3 = \textbf{2} ← \\
Z & \textbf{−3} ↑ & 0 & 0 & 5/2 & 0 & 30 & \\
\bottomrule\noalign{}
\end{tabular}\hfil\null\par\medskip


\endgroup

x₁ enters, s₃ leaves.

\textbf{Tableau 2 (optimal --- no negative entries in Z row)}

\begingroup\small

\par\medskip\noindent\hfil\begin{tabular}{@{}
  >{\raggedright\arraybackslash}p{(\columnwidth - 12\tabcolsep) * \real{0.0857}}
  >{\raggedright\arraybackslash}p{(\columnwidth - 12\tabcolsep) * \real{0.0514}}
  >{\raggedright\arraybackslash}p{(\columnwidth - 12\tabcolsep) * \real{0.0514}}
  >{\raggedright\arraybackslash}p{(\columnwidth - 12\tabcolsep) * \real{0.0514}}
  >{\raggedright\arraybackslash}p{(\columnwidth - 12\tabcolsep) * \real{0.0747}}
  >{\raggedright\arraybackslash}p{(\columnwidth - 12\tabcolsep) * \real{0.0747}}
  >{\raggedright\arraybackslash}p{(\columnwidth - 12\tabcolsep) * \real{0.0805}}@{}}
\toprule\noalign{}
\begin{minipage}[b]{\linewidth}\raggedright
\textbf{Basic}
\end{minipage} & \begin{minipage}[b]{\linewidth}\raggedright
\textbf{x₁}
\end{minipage} & \begin{minipage}[b]{\linewidth}\raggedright
\textbf{x₂}
\end{minipage} & \begin{minipage}[b]{\linewidth}\raggedright
\textbf{s₁}
\end{minipage} & \begin{minipage}[b]{\linewidth}\raggedright
\textbf{s₂}
\end{minipage} & \begin{minipage}[b]{\linewidth}\raggedright
\textbf{s₃}
\end{minipage} & \begin{minipage}[b]{\linewidth}\raggedright
\textbf{RHS}
\end{minipage} \\
\midrule\noalign{}
s₁ & 0 & 0 & 1 & 1/3 & −1/3 & 2 \\
x₂ & 0 & 1 & 0 & 1/2 & 0 & 6 \\
x₁ & 1 & 0 & 0 & −1/3 & 1/3 & 2 \\
Z & 0 & 0 & 0 & 3/2 & 1 & \textbf{36} \\
\bottomrule\noalign{}
\end{tabular}\hfil\null\par\medskip


\endgroup

\textbf{Optimal: x₁ = 2, x₂ = 6, Z = 36.} Shadow prices (dual solution) are read from the Z row under the slacks: y₁ = 0, y₂ = 3/2, y₃ = 1 --- and 4(0) + 12(3/2) + 18(1) = 36 confirms strong duality.

\needspace{130pt}

\hypertarget{p2-01--82-transportation-problem--nwc--modi}{%
\subsection{Transportation Problem --- NWC + MODI}\label{p2-01--82-transportation-problem--nwc--modi}}

\begingroup\small

\par\medskip\noindent\hfil\begin{tabular}{@{}
  >{\raggedright\arraybackslash}p{(\columnwidth - 8\tabcolsep) * \real{0.0920}}
  >{\raggedright\arraybackslash}p{(\columnwidth - 8\tabcolsep) * \real{0.0510}}
  >{\raggedright\arraybackslash}p{(\columnwidth - 8\tabcolsep) * \real{0.0510}}
  >{\raggedright\arraybackslash}p{(\columnwidth - 8\tabcolsep) * \real{0.0510}}
  >{\raggedright\arraybackslash}p{(\columnwidth - 8\tabcolsep) * \real{0.1469}}@{}}
\toprule\noalign{}
\begin{minipage}[b]{\linewidth}\raggedright
\end{minipage} & \begin{minipage}[b]{\linewidth}\raggedright
\textbf{D1}
\end{minipage} & \begin{minipage}[b]{\linewidth}\raggedright
\textbf{D2}
\end{minipage} & \begin{minipage}[b]{\linewidth}\raggedright
\textbf{D3}
\end{minipage} & \begin{minipage}[b]{\linewidth}\raggedright
\textbf{Supply}
\end{minipage} \\
\midrule\noalign{}
S1 & 4 & 6 & 8 & 20 \\
S2 & 5 & 8 & 7 & 30 \\
S3 & 6 & 9 & 5 & 25 \\
Demand & 10 & 25 & 40 & 75 (balanced) \\
\bottomrule\noalign{}
\end{tabular}\hfil\null\par\medskip


\endgroup

\textbf{North-West Corner}: S1D1 = 10, S1D2 = 10, S2D2 = 15, S2D3 = 15, S3D3 = 25 → 5 allocations = m + n − 1 ✓.
Cost = 40 + 60 + 120 + 105 + 125 = \textbf{450}.

\textbf{MODI test} (set u₁ = 0; for basic cells uᵢ + vⱼ = cᵢⱼ): v₁ = 4, v₂ = 6, u₂ = 2, v₃ = 5, u₃ = 0.
Reduced cost of non-basic cell = cᵢⱼ − uᵢ − vⱼ: S1D3 = 3, \textbf{S2D1 = 5 − 2 − 4 = −1} (negative → improve), S3D1 = 2, S3D2 = 3.

Closed loop for S2D1: S2D1(+) → S1D1(−) → S1D2(+) → S2D2(−); θ = min(10, 15) = 10.
New allocation: S1D2 = 20, S2D1 = 10, S2D2 = 5, S2D3 = 15, S3D3 = 25 → cost = 120 + 50 + 40 + 105 + 125 = \textbf{440}.
Recomputing u, v gives all reduced costs ≥ 0 → \textbf{optimal cost 440}.

\needspace{130pt}

\hypertarget{p2-01--83-assignment-problem--hungarian-method}{%
\subsection{Assignment Problem --- Hungarian Method}\label{p2-01--83-assignment-problem--hungarian-method}}

\begingroup\small

\par\medskip\noindent\hfil\begin{tabular}{@{}
  >{\raggedright\arraybackslash}p{(\columnwidth - 6\tabcolsep) * \real{0.0357}}
  >{\raggedright\arraybackslash}p{(\columnwidth - 6\tabcolsep) * \real{0.0492}}
  >{\raggedright\arraybackslash}p{(\columnwidth - 6\tabcolsep) * \real{0.0492}}
  >{\raggedright\arraybackslash}p{(\columnwidth - 6\tabcolsep) * \real{0.0492}}@{}}
\toprule\noalign{}
\begin{minipage}[b]{\linewidth}\raggedright
\end{minipage} & \begin{minipage}[b]{\linewidth}\raggedright
\textbf{J1}
\end{minipage} & \begin{minipage}[b]{\linewidth}\raggedright
\textbf{J2}
\end{minipage} & \begin{minipage}[b]{\linewidth}\raggedright
\textbf{J3}
\end{minipage} \\
\midrule\noalign{}
A & 9 & 2 & 7 \\
B & 6 & 4 & 3 \\
C & 5 & 8 & 1 \\
\bottomrule\noalign{}
\end{tabular}\hfil\null\par\medskip


\endgroup

\[
\underbrace{\begin{pmatrix} 9 & 2 & 7\\ 6 & 4 & 3\\ 5 & 8 & 1 \end{pmatrix}}_{\text{cost}}
\ \xrightarrow{\ \text{rows } -2,-3,-1\ }\
\begin{pmatrix} 7 & 0 & 5\\ 3 & 1 & 0\\ 4 & 7 & 0 \end{pmatrix}
\ \xrightarrow{\ \text{columns } -3,0,0\ }\
\begin{pmatrix} 4 & \boxed{0} & 5\\ \boxed{0} & 1 & 0\\ 1 & 7 & \boxed{0} \end{pmatrix}
\]
\begin{center}\small
Three lines cover all zeros ($= n$), so an optimal assignment exists among the zeros:
A$\to$J2, B$\to$J1, C$\to$J3, with minimum cost $2 + 6 + 1 = 9$.
\end{center}


\needspace{95pt}

\hypertarget{p2-01--84-groups--cayley-tables}{%
\subsection{Groups --- Cayley Tables}\label{p2-01--84-groups--cayley-tables}}

\begin{center}
\renewcommand{\arraystretch}{1.25}
\begin{tabular}{c|cccc}
$+_4$ & 0 & 1 & 2 & 3 \\ \hline
0 & 0 & 1 & 2 & 3 \\
1 & 1 & 2 & 3 & 0 \\
2 & 2 & 3 & 0 & 1 \\
3 & 3 & 0 & 1 & 2 \\
\end{tabular}
\qquad\qquad
\begin{tabular}{c|cccc}
$\cdot$ & $e$ & $a$ & $b$ & $c$ \\ \hline
$e$ & $e$ & $a$ & $b$ & $c$ \\
$a$ & $a$ & $e$ & $c$ & $b$ \\
$b$ & $b$ & $c$ & $e$ & $a$ \\
$c$ & $c$ & $b$ & $a$ & $e$ \\
\end{tabular}

\smallskip
{\small $(\mathbb{Z}_4, +_4)$: cyclic, generated by 1 or 3 \qquad Klein four-group $V_4$: not cyclic; every element has order 2}
\end{center}


\begin{itemize}
\tightlist
\item
  Both have order 4 and are abelian, but are \textbf{not isomorphic} (Z₄ has an element of order 4; V₄ does not).
\item
  Subgroups of Z₄: \{0\}, \{0, 2\}, Z₄. Cosets of H = \{0, 2\}: H and \{1, 3\} → Z₄/H ≅ Z₂.
\item
  Every group of order 4 is isomorphic to Z₄ or V₄; every group of order 6 is Z₆ or S₃.
\end{itemize}

\needspace{125pt}

\hypertarget{p2-01--85-solving-a-recurrence}{%
\subsection{Solving a Recurrence}\label{p2-01--85-solving-a-recurrence}}

aₙ = 5aₙ₋₁ − 6aₙ₋₂, a₀ = 1, a₁ = 0.

\begin{formulas}
\begin{align*}
a_n &= 5a_{n-1} - 6a_{n-2},\qquad a_0 = 1,\ a_1 = 0\\
r^2 - 5r + 6 &= 0 \;\Rightarrow\; r = 2,\ 3\\
a_n &= \alpha\,2^n + \beta\,3^n\\
\alpha + \beta = 1,\quad 2\alpha + 3\beta = 0 &\;\Rightarrow\; \alpha = 3,\ \beta = -2\\
a_n &= 3\cdot 2^n - 2\cdot 3^n \qquad\bigl(\text{check: } a_2 = 12 - 18 = -6 = 5\cdot 0 - 6\cdot 1\bigr)
\end{align*}
\end{formulas}


Non-homogeneous: aₙ = 2aₙ₋₁ + 1 (Hanoi) → particular solution constant c = 2c + 1 ⇒ c = −1; aₙ = A·2ⁿ − 1; a₁ = 1 ⇒ aₙ = 2ⁿ − 1.

\textbf{Generating functions}: the sequence 1, 1, 1, \ldots{} ↔ 1/(1 − x); aₙ = C(n + k − 1, k − 1) ↔ 1/(1 − x)ᵏ. Used to count integer solutions and to solve recurrences.

\needspace{95pt}

\hypertarget{p2-01--86-graph-isomorphism--colouring-checks}{%
\subsection{Graph Isomorphism \& Colouring Checks}\label{p2-01--86-graph-isomorphism--colouring-checks}}

\begin{figure}[!htbp]
\centering
\begin{adjustbox}{max width=\linewidth}
\begin{tikzpicture}[v/.style={circ, minimum size=6.5mm, font=\footnotesize}]
  \begin{scope}
    \node[v, fill=fwine!25] (a) at (0,2) {a}; \node[v] (b) at (2,2) {b}; \node[v, fill=fwine!25] (c) at (2,0) {c}; \node[v] (d) at (0,0) {d};
    \draw[thinline] (a) -- (b) -- (c) -- (d) -- (a) (a) -- (c);
    \node[capt] at (1,-0.8) {$G_1$};
  \end{scope}
  \begin{scope}[shift={(5.5,0)}]
    \node[v] (p) at (1,2.2) {p}; \node[v, fill=fwine!25] (q) at (2.2,1) {q}; \node[v] (r) at (1,-0.2) {r}; \node[v, fill=fwine!25] (s) at (-0.2,1) {s};
    \draw[thinline] (p) -- (q) -- (r) -- (s) -- (p) (q) -- (s);
    \node[capt] at (1,-0.8) {$G_2$};
  \end{scope}
  \node[lbl, anchor=west, text width=40mm, align=left] at (9.2,1) {$a\mapsto q,\ c\mapsto s$\\$b\mapsto p,\ d\mapsto r$\\[3pt]degree-3 vertices shaded};
\end{tikzpicture}
\end{adjustbox}
\caption{Two drawings of the same graph ($K_4$ minus one edge). Matching degree sequences suggest the
isomorphism; exhibiting the mapping proves it.}
\end{figure}


\begin{itemize}
\tightlist
\item
  Both graphs have 4 vertices, 5 edges and degree sequence (3, 3, 2, 2) --- each is K₄ minus one edge.
\item
  Mapping a→q, c→s, b→p, d→r preserves every edge → \textbf{isomorphic}.
\item
  χ(G1) = 3, because the triangle a--b--c needs three colours.
\end{itemize}

Quick invariant checklist for isomorphism: number of vertices and edges, degree sequence, number of cycles of each length, connectivity, bipartiteness. Matching invariants are necessary, not sufficient --- finish by exhibiting a mapping.

\needspace{125pt}

\hypertarget{p2-01--previous-year-questions-pyq-pattern}{%
\section{Previous Year Questions (PYQ pattern)}\label{p2-01--previous-year-questions-pyq-pattern}}

\textbf{Logic}

\begin{enumerate}
\def\labelenumi{\arabic{enumi}.}
\tightlist
\item
  Which is a tautology? (a) p→(p∨q) (b) (p∨q)→p (c) p∧¬p (d) p→¬p --- \textbf{Ans: (a)}
\item
  The contrapositive of "If it rains, the ground is wet" is: \textbf{"If the ground is not wet, it did not rain."}
\item
  ¬∀x (P(x) → Q(x)) is equivalent to: \textbf{∃x (P(x) ∧ ¬Q(x))}
\item
  "Every student has a computer": with S(x), C(x): \textbf{∀x (S(x) → C(x))}
\item
  Number of Boolean functions of 3 variables: 2⁸ = \textbf{256}
\item
  Which is functionally complete? (a) \{∧,∨\} (b) \{↑\} NAND (c) \{∨,↔\} (d) \{∧,→\}... --- \textbf{NAND}
\item
  From p→q and ¬q we infer ¬p --- this is: \textbf{Modus Tollens}
\end{enumerate}

\textbf{Sets, relations, functions}

\begin{enumerate}
\def\labelenumi{\arabic{enumi}.}
\setcounter{enumi}{7}
\tightlist
\item
  Number of reflexive relations on a set of 4 elements: 2¹² = \textbf{4096}
\item
  Number of symmetric relations on 3 elements: 2⁶ = \textbf{64}
\item
  Number of equivalence relations on \{1,2,3,4\}: \textbf{15}
\item
  Number of onto functions from a 4-set to a 3-set: 3⁴ − 3·2⁴ + 3·1 = 81 − 48 + 3 = \textbf{36}
\item
  Relation "≤" on integers is: \textbf{partial order (in fact total)}
\item
  Which lattice is not distributive? \textbf{Pentagon N₅ / Diamond M₃}
\item
  D₃₀ (divisors of 30) is a Boolean algebra because 30 is: \textbf{square-free} (D₁₂ is not)
\end{enumerate}

\textbf{Counting}

\begin{enumerate}
\def\labelenumi{\arabic{enumi}.}
\setcounter{enumi}{14}
\tightlist
\item
  Minimum students so that at least 3 share birthday month: 2×12 + 1 = \textbf{25}
\item
  Number of derangements of 4 letters: \textbf{9}
\item
  Number of ways to arrange letters of "MISSISSIPPI": 11!/(4!4!2!) = \textbf{34650}
\item
  Number of non-negative integer solutions of x+y+z = 10: C(12,2) = \textbf{66}
\item
  Solution of aₙ = 6aₙ₋₁ − 9aₙ₋₂ has form: \textbf{(α + βn)3ⁿ}
\item
  Number of distinct BSTs with 3 keys: \textbf{5} (Catalan)
\end{enumerate}

\textbf{Probability}

\begin{enumerate}
\def\labelenumi{\arabic{enumi}.}
\setcounter{enumi}{20}
\tightlist
\item
  Two dice thrown; probability sum = 7: \textbf{1/6}
\item
  A box: 3 defective of 10. Probability of picking 2 non-defective without replacement: (7/10)(6/9) = \textbf{7/15}
\item
  Bayes problem: Machines A, B produce 60\%, 40\%; defect rates 2\%, 3\%. P(A \textbar{} defective) = 0.012/0.024 = \textbf{1/2}
\end{enumerate}

\textbf{Groups}

\begin{enumerate}
\def\labelenumi{\arabic{enumi}.}
\setcounter{enumi}{23}
\tightlist
\item
  Which is NOT a group? (a) (Z,+) (b) (Q−\{0\},×) (c) (Z,×) (d) (R,+) --- \textbf{Ans: (c)} (no inverses)
\item
  A group of order 7 is: \textbf{cyclic (prime order)}
\item
  Number of generators of cyclic group of order 12: φ(12) = \textbf{4}
\item
  Possible orders of subgroups of a group of order 15: \textbf{1, 3, 5, 15}
\item
  Zₙ is a field iff: \textbf{n is prime}
\item
  Every finite integral domain is a: \textbf{field}
\end{enumerate}

\textbf{Graphs}

\begin{enumerate}
\def\labelenumi{\arabic{enumi}.}
\setcounter{enumi}{29}
\tightlist
\item
  Sum of degrees of a graph with 15 edges: \textbf{30}
\item
  A connected planar graph has 10 vertices and 15 edges. Number of faces: 15 − 10 + 2 = \textbf{7}
\item
  Max edges in simple planar graph with 8 vertices: 3×8 − 6 = \textbf{18}
\item
  Which graph is non-planar? (a) K₄ (b) K₂,₃ (c) K₃,₃ (d) Q₃ --- \textbf{(c)}
\item
  Kₙ has Euler circuit iff: \textbf{n is odd}
\item
  Chromatic number of a cycle with 7 vertices: \textbf{3}
\item
  Number of spanning trees of K₄: 4² = \textbf{16}
\item
  Graph with degree sequence (3,3,3,1) --- possible? Sum = 10 even → yes; (3,3,3,3,1)? Sum = 13 odd → \textbf{No}
\item
  Chromatic number of bipartite graph: \textbf{2}
\end{enumerate}

\textbf{Optimization}

\begin{enumerate}
\def\labelenumi{\arabic{enumi}.}
\setcounter{enumi}{38}
\tightlist
\item
  In a transportation problem with 4 sources and 5 destinations, a non-degenerate BFS has: \textbf{8} allocations
\item
  The best initial basic feasible solution method: \textbf{Vogel\textquotesingle s approximation}
\item
  Assignment problems are solved by: \textbf{Hungarian method}
\item
  The dual of a primal with 3 variables and 5 constraints has: \textbf{5 variables, 3 constraints}
\item
  In PERT, a=2, m=5, b=14: tₑ = (2+20+14)/6 = \textbf{6}, σ² = (12/6)² = \textbf{4}
\item
  Activities on the critical path have total float: \textbf{zero}
\item
  Simplex leaving variable is chosen by: \textbf{minimum ratio test}
\item
  If the objective function is parallel to a binding constraint: \textbf{multiple optimal solutions}
\end{enumerate}

\textbf{More practice questions}

\begin{enumerate}
\def\labelenumi{\arabic{enumi}.}
\setcounter{enumi}{46}
\tightlist
\item
  In simplex, if all entries in the pivot column are ≤ 0, the LPP has: \textbf{an unbounded solution}
\item
  If an artificial variable remains in the basis at a positive level in the optimal tableau: \textbf{the problem is infeasible}
\item
  Shadow price of a constraint equals the optimal value of the corresponding: \textbf{dual variable}
\item
  A transportation solution with fewer than m + n − 1 positive allocations is: \textbf{degenerate}
\item
  In MODI, a negative reduced cost (cᵢⱼ − uᵢ − vⱼ) for an unoccupied cell means: \textbf{the solution can be improved}
\item
  Unbalanced assignment problems are balanced by adding: \textbf{dummy rows or columns with zero cost}
\item
  Number of subgroups of Z₄: \textbf{3}
\item
  Klein four-group is: \textbf{abelian but not cyclic}
\item
  Solution of aₙ = 4aₙ₋₁ − 4aₙ₋₂ has the form: \textbf{(α + βn)2ⁿ}
\item
  Coefficient of x⁵ in 1/(1 − x)³: C(7, 2) = \textbf{21}
\item
  The number of edges in a graph with degree sequence (3, 3, 2, 2): \textbf{5}
\item
  Which statement is true? (a) every abelian group is cyclic (b) every cyclic group is abelian (c) every group of order 4 is cyclic (d) S₃ is abelian --- \textbf{Ans: (b)}
\item
  Order of element 2 in (Z₆, +₆): \textbf{3}
\item
  Number of perfect matchings in K₄: \textbf{3}
\end{enumerate}

\begin{revbox}

\begin{itemize}
\tightlist
\item
  p→q ≡ ¬p∨q ≡ ¬q→¬p · Relations on n: 2\^{}(n²)
\item
  V − E + F = 2 · E ≤ 3V − 6 · Cayley nⁿ⁻² · Odd-degree vertices even
\item
  Lagrange: \textbar H\textbar{} divides \textbar G\textbar{} · Zₙ field iff n prime · generators = φ(n)
\item
  Transport BFS = m + n − 1 · PERT tₑ = (a+4m+b)/6
\end{itemize}

\end{revbox}
